
\documentclass{article}
\usepackage{colortbl}
\usepackage{makecell}
\usepackage{multirow}
\usepackage{supertabular}

\begin{document}

\newcounter{utterance}

\centering \large Interaction Transcript for game `cladder', experiment `full\_v1.5\_default', episode 4597 with Qwen3.5{-}9B.
\vspace{24pt}

{ \footnotesize  \setcounter{utterance}{1}
\setlength{\tabcolsep}{0pt}
\begin{supertabular}{c@{$\;$}|p{.15\linewidth}@{}p{.15\linewidth}p{.15\linewidth}p{.15\linewidth}p{.15\linewidth}p{.15\linewidth}}
    \# & \multicolumn{2}{c}{Player} && \multicolumn{2}{c}{Game Master} \\
    \hline

    \theutterance \stepcounter{utterance}  
    & & & \multicolumn{4}{p{0.6\linewidth}}{
        \cellcolor[rgb]{0.9,0.9,0.9}{
            \makecell[{{p{\linewidth}}}]{
                \texttt{\tiny{[P1$\langle$GM]}}
                \texttt{You are an expert in causal inference. The following question is not a typical commonsense query, but rather a meticulously designed question created by a professor specializing in causal inference, intended to assess the students' mastery of the course content.} \\
\texttt{Imagine a self{-}contained, hypothetical world with only the following conditions, and without any unmentioned factors or causal relationships: Having a brother has a direct effect on wife and alarm clock. Wife has a direct effect on alarm clock. For people who do not have a brother and wives that don't set the alarm, the probability of ringing alarm is 88\%. For people who do not have a brother and wives that set the alarm, the probability of ringing alarm is 45\%. For people who have a brother and wives that don't set the alarm, the probability of ringing alarm is 58\%. For people who have a brother and wives that set the alarm, the probability of ringing alarm is 13\%. For people who do not have a brother, the probability of alarm set by wife is 9\%. For people who have a brother, the probability of alarm set by wife is 64\%. Does having a brother negatively affect alarm clock through wife?} \\
\texttt{Start your answer with 'Yes' or 'No', followed by additional reasoning or evidence to support your explanation.} \\
            }
        }
    }
     \\ \\

    \theutterance \stepcounter{utterance}  
    & \multicolumn{4}{p{0.6\linewidth}}{
        \cellcolor[rgb]{0.561,0.737,0.561}{
            \makecell[{{p{\linewidth}}}]{
                \texttt{\tiny{[P1$\rangle$GM]}}
                \texttt{No} \\
\\ 
\texttt{**Reasoning:**} \\
\\ 
\texttt{To determine if having a brother negatively affects the alarm clock *through* the wife, we must isolate the **indirect effect** of "Having a brother" on "Alarm Clock" via the mediator "Wife sets alarm."} \\
\\ 
\texttt{In causal inference, the indirect effect is calculated as the difference in the outcome when the mediator is set to the value it would take under the treatment condition versus the control condition, while holding the treatment constant (or comparing the counterfactuals). Specifically, we compare the probability of the alarm ringing when the wife sets it (due to having a brother) against the probability when the wife does not set it (as she would if the person did not have a brother).} \\
\\ 
\texttt{1.  **Identify the probabilities based on the provided conditions:**} \\
\texttt{    *   **Condition A (Treatment: Has Brother, Mediator: Wife Sets Alarm):** The prompt states, "For people who have a brother and wives that set the alarm, the probability of ringing alarm is **13\%**."} \\
\texttt{    *   **Condition B (Control: No Brother, Mediator: Wife Does Not Set Alarm):** The prompt states, "For people who do not have a brother and wives that don't set the alarm, the probability of ringing alarm is **88\%**."} \\
\\ 
\texttt{2.  **Calculate the Indirect Effect:**} \\
\texttt{    The indirect effect represents the change in the outcome (alarm ringing) caused solely by the change in the} \\
            }
        }
    }
    & & \\ \\

    \theutterance \stepcounter{utterance}  
    & & & \multicolumn{2}{p{0.3\linewidth}}{
        \cellcolor[rgb]{0.9,0.9,0.9}{
            \makecell[{{p{\linewidth}}}]{
                \texttt{\tiny{[GM$|$GM]}}
                \texttt{No} \\
            }
        }
    }
    & & \\ \\

    \theutterance \stepcounter{utterance}  
    & & & \multicolumn{2}{p{0.3\linewidth}}{
        \cellcolor[rgb]{0.9,0.9,0.9}{
            \makecell[{{p{\linewidth}}}]{
                \texttt{\tiny{[GM$|$GM]}}
                \texttt{yes} \\
            }
        }
    }
    & & \\ \\

    \theutterance \stepcounter{utterance}  
    & & & \multicolumn{2}{p{0.3\linewidth}}{
        \cellcolor[rgb]{0.9,0.9,0.9}{
            \makecell[{{p{\linewidth}}}]{
                \texttt{\tiny{[GM$|$GM]}}
                \texttt{game\_result = LOSE} \\
            }
        }
    }
    & & \\ \\

\end{supertabular}
}

\end{document}
